\section{Targeted analysis of the recovered local proofs}

The current review includes four complementary October audits:
global/source assembly and a delicate C1 endpoint; the later full R1
proof; noncommuting-residue constrained geometry; and C2 source
dispatch. The separate audit notes give exact paths, equation tags,
hypotheses and limits. Their findings are summarized here without
converting scoped checks into certification of all 215 components.

\subsection{R1: the effective reserve and the whole parent}

The later R1 theorem expressly removes the earlier restriction that
\(jT\) be infinitesimal. Its compact source variables satisfy
\[
 0<p<\min(1,j),\quad 0<H<1,\qquad
 B,T,jB,TB\text{ positive infinitesimals},
\]
while \(j\) and \(jT\) need not be bounded. A literal inverse map to
the original stationary source is printed.

For \(h_0>0\), \(n\ge0\), the central exact completion is
\begin{equation}
 h_0N_0^2+n(N_0+jb)^2
 =(n+h_0)\left(N_0+\frac{nj}{n+h_0}b\right)^2
 +\frac{nh_0j^2}{n+h_0}b^2.
 \label{eq:effective}
\end{equation}
The denominator is strictly positive and pure source.
Writing
\[
 c_b=\frac{nj}{n+h_0},\qquad M_e=\frac{nh_0j^2}{n+h_0},
\]
the usable reserve is \(M_e\), not the uncompleted mass \(nj^2\).
At \(n=0\), both quantities vanish literally. For \(n>0\),
\[
 \frac{c_b^2}{M_e}=\frac{n}{h_0(n+h_0)}\le\frac1{h_0},
\]
which controls the shift against the retained current.

The full-normal proof retains both unrestricted state normals, both
first-normal arrays, and their quadratic cross term. Scalar, radial
and circuit identities buy helper pivots; strict Gram then applies
to the remaining complete expression. The scalar restriction alone
would not suffice. \src{R1N; FN}

The contact rows use
\[
 U=\frac{u}{\sqrt{jB}},\quad V=\frac{v}{\sqrt{TB}},
 \quad Z_*=\frac{z}{\sqrt{jT}},
\]
and the exact circuit
\[
\begin{aligned}
 \Psi&=U^2+V^2+Z_*^2+U^2V^2Z_*^2-4UVZ_*\\
 &=(U-V)^2+Z_*^2(1-UV)^2+2UV(1-Z_*)^2.
\end{aligned}
\]
The rows can be unbounded. The positive state ratio \(z\) remains
arbitrary inside them. The square root of \(UV\) factors into
pure-source and pure-testing quantities.

The opposite second-normal contribution is treated through the whole
squared multiplier
\[
 T(jT)v^2\bigl(\alpha-BTCHu^2\bigr)^2.
\]
The newer cost estimates cancel \(jT\) exactly and retain all higher
terms and the cross term. They do not compare source and testing
infinitesimals. \src{FN24; R1N11}

For outer tests put
\[
 C_t=v(uX-Z),\qquad B_t=X+uZ,\qquad g_*=\frac{CD}{C+D}.
\]
A pure-source harmonic completion gives \(g_*C_t^2\).
The exact relation
\[
 uB_t-C_t/v=(1+u^2)Z
\]
produces the two-row outer Gram when \(v\) is noninfinitesimal.
The small-\(v\) four-row identity and the remaining six-row identity
retain every state coordinate. Real cap choices use positivity on
event kernels; excess coefficients stay in whole event squares.

The audit directly checked these outer identities, including the
zero row relation used in the six-row adjustment. The state/sign
join covers \(wz\le0\) with a pure-testing nonnegative core, and
normalizes \(w=1\) when \(wz>0\) while keeping \(z>0\) arbitrary.
Homogeneous formulas include literal endpoints, and complex states
split into real and imaginary parts. \src{R1N12--R1N24; SJ}

\textbf{Finding.}
The later R1 argument is present and substantive. No concrete gap
was found in the examined allocation and joins. Older records marking
R1 open should not be substituted for this later body.
The large source-polynomial, scalar-circuit and radial coefficient
expansions were not all independently recomputed in this review.

\subsection{A shorter R1 reserve formulation}

The later proof map combines several source cases using
\[
 \rho=Bj^2,\quad\delta=1-p,\quad S=\delta+\rho+M_e,
 \qquad p^2\le10M_e+4\rho,\quad \frac1{10}\le S<9.
\]
Weighting the exact budgets by \(\delta/S,\rho/S,M_e/S\)
avoids isolated inverse openings. \(S\) is a positive pure-source
unit, and \(M_e=0\) is retained without division by \(M_e\).
This is a useful shorter exposition, retaining the full parent,
circuit, radial and normal identities as inputs. \src{RM}

\subsection{The C1 JC4/SN endpoint}

The reviewed endpoint has \(S,\kappa,\nu\) infinitesimal,
\(\ell=1-\nu\), \(q=\nu/\ell\), and \(n_s=\beta/q^2\) bounded.
The exact source identities give
\[
 b_0^2=O(\kappa\theta n_s),\qquad
 \frac{D}{Sq^2}\longrightarrow\frac12,\qquad
 \frac e\kappa\longrightarrow1,\qquad
 \frac h\nu\longrightarrow1.
\]
This retains literal \(\beta=0\) and excludes \(D=0\).
No ordering among the independent infinitesimal scales is assumed.
The local \(e=xh/\sigma=\kappa h/\nu\) is explicitly defined and is
not the earlier native parameter bearing the same letter.

The complete local parent keeps unrestricted state normals
\(X_e,N_y\). Its exact relation
\[
 R=\frac{\kappa\sqrt S\,G-uN}{\ell}
\]
puts every \(R\)-term into an infinitesimal Gram perturbation.
The helper identity has the usable form
\[
 \mathrm{Cost}=\ell^2(wB_u)^2+(\kappa\delta W)^2,
\]
where the cost's coefficients on strict base rows have zero residue.
It buys any fixed positive helper pivot on \(wB_u\), together with
a genuine square, at infinitesimal cost. The analogous \(wB\) pivot
handles the remaining bounded cross terms.
When \(n_s=0\), its zero row is omitted.

The remote coverage has three genuine cases.
Noninfinitesimal Alice tests use complete four-state Grams since
\(C,A\) are unbounded. Infinitesimal Alice with unmatched Bob uses
a homogeneous six-root Gram with scalar-block determinant
\[
 \frac14(\bar s^2-\bar\ell)^2.
\]
On the SN source \(\bar\ell=1\); the remaining matched Bob tests have
\(c\) infinitesimal and \(s\) a unit, so \(w=c/s\) is infinitesimal,
exactly the local SN domain. State ratios remain in roots throughout.

\textbf{Finding.}
The source transition, full-state allocation and local/remote union
are coherent on the examined formulas. Some coefficient expansions
and the other C1 endpoint cells were not independently rederived.
\src{Current assembly audit}

\subsection{Noncommuting residue tests}

The A6 route uses the geometry of zero events at an entangled residue:
the supported zero events form a matching in the event eight-cycle.
At most three are independent state-normal coordinates. A fourth
matching event enters a separate terminal.

At each source scale, impose the current hard event amplitudes
\emph{exactly} zero in the coefficient field, divide by the positive
pure-source next scale \(\sigma\), and then take residue.
Physicality makes the resulting constrained residue nonnegative.
This does not delete actual terms from the final tensor expression.
The weighted-contact lemma returns to the full parent with all normal
deviations retained.

If the signed coefficient becomes visible, the one-, two- or
three-hard-event constrained curvature argument applies. Otherwise
another matching edge is added. The process terminates because
there are at most four matching edges, not because an arbitrary
infinitesimal series is assumed to stabilize.

The focused audit checked all eight one-hard derivative coefficients
and all 24 two-hard activation conditions from the printed determinant
dictionary, the explicit negative witness states and the nonzero
signed-core differential on the residual matching.
The weighted and unequal-scale lemmas use only real residue Jacobian
inversion and source-field implicit-function transfer; they do not
invert a mixed tensor.

\textbf{Finding.}
No defect was found in the reviewed geometry, constrained curvature
and induction interface. The original full determinant identity and
the complete arbitrary-scale four-matching terminal remain named
mathematical inputs. A6 does not cover commuting-angle residues.
\src{Current A6 audit}

\subsection{C2 source dispatch and boundary logic}

The HJ source partition reaches ten rows R1--R10 or four regular exits.
The two excluded intersections have source-equation justifications:
\[
 \text{all six weights bounded}
 \ \Longleftrightarrow\
 t,h,\ell,j,r,g\text{ are units},
\]
contradicting the weak total scale, and the weak-\(C\) criterion
excludes exactly the simultaneous-small \(r,j\) intersection.

The formal C1 comparison has
\[
 \beta_1=-\frac{r\delta n}{eP_n}.
\]
At \(n=0\) it is an actual C1 boundary source and the precise local
C1 theorem is used. At \(n>0\), the comparison has negative
\(\beta_1\); the proof instead checks finite coefficient identities.
It does not invoke physical C1 positivity for this formal source.

The complete CT reflection preserves the strict native gaps. For example,
\[
 g'=\frac{a^2j}{b},\qquad j'=\frac gb,
\]
\[
 1-r'=\frac{b-a^2zr}{b},\qquad
 1-z'h'=\frac{a^2t}{b},\qquad
 b'-a^2z'r'=\frac{a^2(1-r)}b.
\]
All are positive by the original source domain. Both \(n=0\) and the
total-scale predicate are preserved.

The nine ordered angle types and complementary-event cases are
exhaustive. The single-angle argument explicitly includes additional
opening extensions beyond its earliest theorem header.
The inspected dispatch does not call the entire C1 theorem in a circle;
its C1 imports are named local bodies.

\textbf{Finding.}
No uncovered source cell, false reflection, lost \(n=0\) boundary or
whole-theorem routing cycle was found. Native entry, positive-scale
coverage, A6 and the detailed local conductors remain substantive
inputs; dispatch verification is not independent rederivation of
every leaf. \src{Current C2 dispatch audit}

